\subsection{Three source corrections received from the Algebra review}\label{three-source-corrections-received-from-the-algebra-review}

These arguments concern three identified findings in two other chapters. They do not represent a complete review of either chapter. The proposed source edits consist of one vector index, two tensor subscripts, and one pair of parentheses. Complete arguments and their propagation remain identified editorial material.

The original source is by the Stacks Project authors, at commit a04446e57ec1fbc252a871afcec7752fb2807b14. The earlier receiving records are ALGEBRA-RECON-1069, ALGEBRA-RECON-997, and ALGEBRA-RECON-1333, respectively. This note supplies the exact receiving calculations. No novelty, independent human review, or formal proof certification is claimed.

\subsubsection{CC-FITTING-INDEX: the appended kernel vectors}\label{cc-fitting-index-the-appended-kernel-vectors}

Source: \href{https://github.com/stacks/stacks-project/blob/a04446e57ec1fbc252a871afcec7752fb2807b14/more-algebra.tex\#L1355}{More on Algebra, Fitting ideals, lines 1355--1411}. The erroneous index is at line 1385.

Keep the source ring \(R\), its finite module \(M\), a surjection \(\pi:R^{\oplus n}\to M\), and its kernel \(K\). Fix \(0\leq k\leq n\) and put \(q=n-k\), only as an additional name for that integer. For chosen vectors \(z_1,\ldots,z_{n-k}\in K\), retain their full coordinate matrix \[
 Z=(z_{ij})_{1\leq i\leq n,\ 1\leq j\leq n-k}.
\] After adding \(n'\) generators that map to zero, the surjection is \[
 \pi':R^{\oplus n}\oplus R^{\oplus n'}\longrightarrow M,
 \qquad (a,b)\longmapsto\pi(a),
 \qquad K'=K\oplus R^{\oplus n'}.
\] The required list has exactly \(n+n'-k\) vectors: \[
 z'_j=(z_j,0)\quad(1\leq j\leq n-k),\qquad
 z'_{n-k+j'}=(0,e_{j'})\quad(1\leq j'\leq n'),
\tag{F1}
\] where \(e_{j'}\) is the \(j'\)-th standard coordinate vector of \(R^{\oplus n'}\). Its coordinate matrix is precisely \[
 Z'=
 \begin{pmatrix}
  Z&0_{n\times n'}\\
  0_{n'\times(n-k)}&1_{n'}
 \end{pmatrix}.
\tag{F2}
\] The printed index \(n+j'\) leaves the positions \(n-k+1,\ldots,n\) undefined when \(k>0\), and its largest indices exceed the stated list. For example \(n=2,k=1,n'=1\) requires two columns numbered 1 and 2, whereas the printed rule calls the added column 3. Equation (F1) repairs exactly this index; neither the vector nor its coordinates changes.

We prove the ideal comparison used by the source. Every \((n+n'-k)\)-row minor of (F2) uses all its columns. A selected row set that omits a bottom row makes the corresponding added coordinate column zero, so its determinant is zero. A selected row set containing every bottom row has \(n-k\) top rows. With both rows and columns in their original increasing order, its determinant equals the corresponding \((n-k)\)-row minor of \(Z\): its matrix is block diagonal with bottom block \(1_{n'}\). Every such minor of \(Z\) occurs. This proves equality of these two minor ideals, including \(n'=0\) and \(k=n\). In the latter case the empty top determinant is 1, and the bottom identity matrix provides the unit minor.

Conversely choose any \(n+n'-k\) vectors in \(K'\), with coordinate matrix \(W\), and project each of them to \(K\). Write their projected matrix \(P\) and their bottom coordinate matrix \(Q\), so \[
 W=\begin{pmatrix}P\\Q\end{pmatrix},
 \qquad P\in\operatorname{Mat}_{n\times(n+n'-k)}(R),
 \quad Q\in\operatorname{Mat}_{n'\times(n+n'-k)}(R).
\] For a selected \((n+n'-k)\)-row minor, let \(s\) be the number of selected bottom rows. Expansion along these \(s\) rows is the sum over all \(s\)-element column subsets of the selected bottom determinant times the complementary \((n+n'-k-s)\)-row determinant of \(P\), with its row-selection and column-selection sign retained. Since \(s\leq n'\), the latter determinant has size at least \(n-k\). Expanding it along any \(n-k\) selected rows writes it as a signed sum of \((n-k)\)-row minors of \(P\) multiplied by complementary minors. Each of its columns is one of the original projected vectors in \(K\). Every term therefore belongs to the source ideal generated by the \((n-k)\)-row minors of kernel vectors. For \(k=n\) this containment is containment in the unit ideal. This proves the other inclusion without discarding any sign or coordinate contribution.

For completeness, this kernel-vector description agrees with the presentation matrix description even if the relation index set is infinite: every kernel vector is a finite linear combination of its columns. Multilinearity expands a chosen determinant into a finite sum of coefficient products times minors of the presentation matrix. Terms with repeated columns vanish by alternation. Conversely, each presentation column is itself a kernel vector, so every presentation minor occurs in the kernel-vector ideal.

The comparison between different generating surjections has an explicit ambient isomorphism. For \(\rho:R^{\oplus m}\to M\), choose lifts \[
 U:R^{\oplus m}\to R^{\oplus n},\quad \pi U=\rho,\qquad
 V:R^{\oplus n}\to R^{\oplus m},\quad \rho V=\pi.
\] These exist by choosing a lift of the image of each standard basis vector. On \(R^{\oplus n}\oplus R^{\oplus m}\), the automorphisms \[
 T_\pi(a,b)=(a+Ub,b),\quad T_\pi^{-1}(a,b)=(a-Ub,b),
\] \[
 T_\rho(a,b)=(a,b+Va),\quad T_\rho^{-1}(a,b)=(a,b-Va)
\] satisfy \[
 (\pi,0)T_\pi=(\pi,\rho)=(0,\rho)T_\rho.
\] Thus \(T_\rho T_\pi^{-1}\) identifies \(\ker(\pi)\oplus R^{\oplus m}\) with \(R^{\oplus n}\oplus\ker(\rho)\), and both composites with the surjections are the displayed maps to the original \(M\). The minors of a matrix multiplied by an ambient automorphism belong to the original minor ideal by the determinant expansion. Multiplication by the inverse proves the reverse inclusion. A permutation of the two coordinate blocks changes only the corresponding determinant signs, which are units and remain present in the expansion. This proves the independence comparison with the corrected indices.

For \(k>n\), the \(k\)-th Fitting ideal for a presentation with \(n\) generators is the unit ideal. This convention is consistent with stabilization: if \(0<n+n'-k\leq n'\), select that many added coordinate vectors and the same bottom coordinates to obtain determinant 1; if \(n+n'-k\leq0\), the unit convention applies directly. No claim about a negative number of chosen kernel vectors is needed.

Propagation: the Jacobian-defect calculations in the receiving Algebra note use presentation independence. Their matrices and Fitting ideals are unchanged; equation (F1) gives the correctly indexed comparison. The source theorem is retained, and the longer verification stays in this editorial note.

\subsubsection{\texorpdfstring{CC-CONORMAL-BASE: the tensor product is over \(B\)}{CC-CONORMAL-BASE: the tensor product is over B}}\label{cc-conormal-base-the-tensor-product-is-over-b}

Source: \href{https://github.com/stacks/stacks-project/blob/a04446e57ec1fbc252a871afcec7752fb2807b14/more-algebra.tex\#L8458}{More on Algebra, the transitivity lemma, lines 8458--8505}. The wrong tensor subscripts occur at lines 8502 and 8504.

Retain all the source maps \[
 A\longrightarrow B\longrightarrow C,\qquad
 P=A[x_s\mid s\in S]\twoheadrightarrow B=P/I,\qquad
 B[y_1,\ldots,y_m]\twoheadrightarrow C.
\] Put \(T=P[y_1,\ldots,y_m]\), let \(K=\ker(T\to C)\), and let \(J=\ker(B[y_1,\ldots,y_m]\to C)=K/IT\). Here \(IK\) means the ideal \((IT)K\) of \(T\), as in the source. The \(B\)-action on \(I/I^2\) is induced by the original quotient \(P\to P/I\); the \(B\)-action on \(C\) is the original given map.

For these data, without an additional regularity assumption, the map \[
 \Phi:(I/I^2)\otimes_B C\xrightarrow{\ \sim\ }IT/IK,
 \qquad (i\bmod I^2)\otimes(t\bmod K)\longmapsto it\bmod IK
\tag{C1}
\] is an isomorphism of \(C\)-modules. If \(i\) changes by an element of \(I^2\), its product with \(t\) changes by an element of \(I^2T\subset IK\), since \(IT\subset K\). If \(t\) changes by an element of \(K\), its product with \(i\) changes by an element of \(IK\). For a lift \(p\in P\) of \(b\in B\), the two balanced tensors \((ip\bmod I^2)\otimes\bar t\) and \((i\bmod I^2)\otimes\overline{pt}\) have the same image. A different lift of \(b\) makes their products differ by an element of \(I^2T\). This proves every required well-definedness and balance relation.

Here is an inverse retaining every polynomial coefficient. Each \(q\in IT\) has a unique finite polynomial expansion \[
 q=\sum_{\nu\in\mathbf N^m}i_\nu y^\nu,\qquad i_\nu\in I,
\] and define \[
 \Psi(q\bmod IK)=
 \sum_{\nu\in\mathbf N^m}(i_\nu\bmod I^2)\otimes\overline{y^\nu}.
\tag{C2}
\] The sum is finite. Coefficient multiplication by \(p\in P\) is balanced with its original image in \(B\); multiplication by \(y_t\) is balanced with its image in \(C\). Consequently the expression for \(q=i\,h\), with \(i\in I,h\in T\), is \((i\bmod I^2)\otimes\bar h\). Every element of \(IK\) is a finite sum of terms \(i\,h\,k\), where \(k\in K\); each maps to \((i\bmod I^2)\otimes\overline{hk}=0\). Thus (C2) descends to the quotient. Applying (C1) to (C2) returns the complete polynomial \(q\) modulo \(IK\). Applying (C2) to (C1) on an elementary tensor returns that tensor by expanding \(t\) into all its polynomial coefficients. Both composites are identities, and the displayed balance proves \(C\)-linearity.

The source additionally assumes \(B\to C\) is a local complete intersection. We verify the exact consequence it needs: \[
 IT\cap K^2=IK.
\tag{C3}
\] The source definition makes \(J\) a Koszul-regular ideal, hence locally generated by a sequence whose first Koszul homology vanishes. For the relevant definition and original proofs see \href{https://github.com/stacks/stacks-project/blob/a04446e57ec1fbc252a871afcec7752fb2807b14/more-algebra.tex\#L8100}{regular ideals, lines 8100--8137}, \href{https://github.com/stacks/stacks-project/blob/a04446e57ec1fbc252a871afcec7752fb2807b14/more-algebra.tex\#L7471}{the intersection calculation, lines 7471--7520}, and \href{https://github.com/stacks/stacks-project/blob/a04446e57ec1fbc252a871afcec7752fb2807b14/more-algebra.tex\#L8234}{its ideal form, lines 8234--8244}.

The needed calculation is as follows. Let \(D\) be a ring, \(L\) an ideal, and \(g_1,\ldots,g_r\in D\) whose images in \(D/L\) have vanishing first Koszul homology. For \(h=\sum_i g_i a_i\in L\), the vector \(\sum_i\bar a_i e_i\) is a Koszul cycle. With the convention \[
 d_2(e_i\wedge e_j)=\bar g_i e_j-\bar g_j e_i\quad(i<j),
\] there are coefficients \(\bar c_{ij}\in D/L\) such that it equals \(\sum_{i<j}\bar c_{ij}(\bar g_i e_j-\bar g_j e_i)\). Choose lifts \(c_{ij}\in D\) and set \[
 \ell_i=a_i+\sum_{j>i}c_{ij}g_j-\sum_{h<i}c_{hi}g_h\in L.
\] Then, keeping both signed terms for every pair, \[
 \sum_i g_i a_i
 =\sum_i g_i\ell_i+
   \sum_{i<j}(-c_{ij}g_i g_j+c_{ij}g_j g_i)
 =\sum_i g_i\ell_i.
\] This proves \(L\cap(g_1,\ldots,g_r)=L(g_1,\ldots,g_r)\); the other containment follows by multiplication. If \(Q=L+(g_1,\ldots,g_r)\), an element of \(L\cap Q^2\) has the form \[
 h=\sum_{i,j}b_{ij}g_i g_j+\sum_i\ell'_i g_i+\ell',
 \quad b_{ij}\in D,\quad\ell'_i\in L,\quad\ell'\in L^2.
\] The first sum belongs to \(L\cap(g_1,\ldots,g_r)\), and is therefore in \(L(g_1,\ldots,g_r)\). Every displayed summand lies in \(LQ\). Hence \(L\cap Q^2=LQ\), including the empty-sequence case.

Apply this with the actual localized ring \(D=T_g\), \(L=(IT)_g\), and \(Q=K_g\) around each prime of \(T\) containing \(K\). Such a \(g\) is the lift of a source neighborhood on which \(J\) has Koszul-regular generators. At a prime not containing \(K\), \(K\) localizes to the unit ideal and both sides of (C3) localize to \(IT\). Localization commutes with these products and this finite intersection: an element in both localized ideals has two denominator witnesses, and their product is a common witness. Thus the quotient \((IT\cap K^2)/IK\) vanishes at every prime. If this module had a nonzero element, a maximal ideal containing its annihilator would give a localization in which it survives. This proves (C3) globally. No finite generation of \(I\), \(K\), or that quotient was assumed.

By (C1) and (C3), the natural map \[
 (I/I^2)\otimes_B C\longrightarrow K/K^2,\qquad
 \bar i\otimes\bar t\longmapsto it\bmod K^2
\] is injective. The map \(K/K^2\to J/J^2\) is surjective, and its kernel is \((IT+K^2)/K^2\), identified with \(IT/(IT\cap K^2)\). Therefore the exact row in the source is precisely \[
 0\longrightarrow (I/I^2)\otimes_B C
 \longrightarrow K/K^2\longrightarrow J/J^2\longrightarrow0.
\tag{C4}
\] The polynomial differential row retains the separate \(dx_s\) and \(dy_t\) coordinates: \[
 0\longrightarrow
 \bigoplus_{s\in S}C\,dx_s
 \longrightarrow
 \left(\bigoplus_{s\in S}C\,dx_s\right)\oplus
 \left(\bigoplus_{t=1}^{m}C\,dy_t\right)
 \longrightarrow
 \bigoplus_{t=1}^{m}C\,dy_t
 \longrightarrow0.
\tag{C5}
\] The conormal differentials send an equation to its original universal derivative tensored with \(C\). For \(i\in I,t\in T\), the identity \(d(it)=t\,di+i\,dt\) becomes \(t\,di\), because \(i\) maps to zero in \(C\). This proves commutation at the left; reducing an equation modulo \(I\) and taking its relative \(B\)-derivative proves commutation at the right.

Consequently the source's two-term-complex comparison and six-term sequence have exactly the required conormal row (C4). More explicitly, the connecting map takes a class in the kernel of the conormal differential for \(J\), lifts its equation to \(K\), differentiates in \(T/A\), and retains the \(dx_s\) component. Its \(dy_t\) component is zero by the cycle condition. Its class in \(\Omega_{B/A}\otimes_B C\) is unchanged by changing the lift: the difference lies in \(IT\), and its derivative is a boundary from the left member of (C4). A change by an element of \(J^2\) can be lifted to \(K^2\), whose derivative vanishes after tensoring with \(C\). The resulting kernel/image equalities follow directly: a middle cycle mapping to zero in the right kernel is a left cycle by (C4) and (C5); a right cycle whose connecting class is zero can have its lift corrected by a left conormal element to become a middle cycle; and a left differential class dying in the middle cokernel is the connecting image of the right image of an element whose middle boundary represents it. Exactness at the other cokernel positions is right exactness of (C4) and (C5). The initial kernel map is injective because the initial conormal map is injective. These are the original source maps, with the tensor over \(B\).

The wrong subscript has an explicit counterexample satisfying the source hypothesis. Let \[
 A=k,\quad P=k[x],\quad I=K=(x^2),\quad
 B=C=k[x]/(x^2),
\] and take the identity \(B\to C\), presented with no \(y\)-variables. Its kernel \(J=0\) is generated by the empty Koszul-regular sequence, so this map is a local complete intersection. The source quotient is \[
 IT/IK=(x^2)/(x^4)=I/I^2,
\] with \(k\)-basis \(u=x^2\bmod x^4,\ v=x^3\bmod x^4\). In contrast, \((I/I^2)\otimes_k C\) has the four-element basis \[
 u\otimes1,\quad u\otimes\bar x,\quad v\otimes1,\quad
 v\otimes\bar x.
\] Its multiplication map to \(I/I^2\) sends these four basis vectors to \(u,v,v,0\), respectively. Its kernel has the independent basis \[
 u\otimes\bar x-v\otimes1,\qquad v\otimes\bar x.
\] Indeed a vector with coefficients \(a,b,c,d\) maps to \(a\,u+(b+c)v\), so it is in the kernel exactly when \(a=0,c=-b\). This proves the erroneous tensor cannot be identified with the source quotient and cannot supply the asserted injection. The \(B\)-balanced tensor identifies the middle two basis vectors and makes the last zero, exactly as (C1) requires.

Propagation: replace the two printed \(A\) subscripts by \(B\). The diagram, the six-term statement, and the existing Algebra cotangent calculations already use the correct \(B\)-action. The exact identification (C1) needs no local-complete-intersection assumption; the injection into \(K/K^2\) uses the explicitly proved intersection equality. This distinction is a proved dependency fact, not a proposed new theorem attributed to the source authors.

\subsubsection{\texorpdfstring{CC-SERIES-FACTOR: retain the coefficient's factor of \(x^{-1}\)}{CC-SERIES-FACTOR: retain the coefficient's factor of x\^{}\{-1\}}}\label{cc-series-factor-retain-the-coefficients-factor-of-x-1}

Source: \href{https://github.com/stacks/stacks-project/blob/a04446e57ec1fbc252a871afcec7752fb2807b14/examples.tex\#L1428}{Examples, the original series and the localized comparison, lines 1428--1483}. The faulty recurrence is at line 1481.

Keep the source field \(k\), the original indeterminate \(x\), and all coefficients of its stipulated transcendental series \[
 z=\sum_{i=1}^{\infty}a_i x^i\in k[[x]],\qquad
 z_j=x^{-j}\left(z-\sum_{i=1}^{j-1}a_i x^i\right)
     =\sum_{i=j}^{\infty}a_i x^{i-j}\quad(j\geq1).
\] The coefficient calculation in the original Laurent series field gives \[
 xz_{j+1}
 =\sum_{i=j+1}^{\infty}a_i x^{i-j}
 =z_j-a_j,\qquad
 z_{j+1}=x^{-1}(z_j-a_j).
\tag{S1}
\] The printed expression \(x^{-1}z_j-a_j\) minus the right side of (S1) is \(a_j(x^{-1}-1)\). The series is transcendental over \(k(x)\), so at least one of its coefficients is nonzero. At any such index this difference is nonzero in \(k((x))\). Neither the coefficient nor the power of \(x\) can be discarded.

Let \(R=k[x,z_1,z_2,\ldots]\subset k[[x]]\), as in the source. Its Laurent localization satisfies the exact equality of subrings of \(k((x))\) \[
 R[1/x]=k[x,x^{-1},z].
\tag{S2}
\] The inclusion from right to left uses \(z=xz_1\in R\) and the adjoined inverse of \(x\). For the other inclusion, the full formula for every \(z_j\) displays it in the right-hand ring. Thus both maps are induced by the same original elements in \(k((x))\).

We also verify the residue ideal used by the localized comparison. Each \(z_j\) is transcendental over \(k(x)\): an algebraic relation for \(z_j\), combined with \(z=x^jz_j+\sum_{i<j}a_i x^i\), would make \(z\) algebraic. Hence \(R_j=k[x,z_j]\) is the polynomial ring on the two displayed elements. The exact transition maps \[
 R_j\longrightarrow R_{j+1},\qquad
 x\longmapsto x,\quad z_j\longmapsto xz_{j+1}+a_j
\] are inclusions in \(k[[x]]\), and their union is \(R\). Modulo \(x-1\) they become the polynomial-ring isomorphisms \[
 k[z_j]\longrightarrow k[z_{j+1}],\qquad
 z_j\longmapsto z_{j+1}+a_j.
\] Consequently \[
 R/(x-1)\xrightarrow{\sim}k[z],\qquad
 x\longmapsto1,\quad
 z_j\longmapsto z-\sum_{i=1}^{j-1}a_i.
\tag{S3}
\] This can be checked without assuming quotient-commutation: define the compatible maps on each \(R_j\) by these formulas; they induce the map on the union \(R\). Its inverse sends the polynomial indeterminate \(z\) to the class of the original \(xz_1\). The transition relations verify both composites on every generator. Thus its kernel is exactly \((x-1)\).

The source ideal \[
 \mathfrak n=(x-1,z_1,z_2+a_1,z_3+a_1+a_2,\ldots)
\] is therefore \((x-1,z)\), with residue field \(k\). In particular \(x\notin\mathfrak n\). Equality (S2) gives the actual comparison \[
 k[x,x^{-1},z]_{(x-1,z)}
 \xrightarrow{\sim} R_{\mathfrak n}.
\tag{S4}
\] The forward map sends the original \(x,x^{-1},z\) to those same elements. For its inverse, send each \(z_j\) to \[
 x^{-j}\left(z-\sum_{i=1}^{j-1}a_i x^i\right).
\] These formulas respect every ring relation because (S2) identifies the two rings inside \(k((x))\). An element of \(R\setminus\mathfrak n\) maps to an element whose value at \(x=1,z=0\) is nonzero by (S3), and hence is invertible in the left side of (S4). Conversely a denominator outside \((x-1,z)\) maps outside the maximal ideal on the right. Both maps therefore extend to the indicated localizations, and their composites are the identity on all numerators and denominators. This proves the source's localized comparison with every original coefficient retained.

Propagation: the correction changes the recurrence used to display the inverse in (S4), while the original ring, prime ideal and localization map remain the same. The receiving Algebra descent note already uses (S1); its finite étale construction does not require a new source object. Its complete argument remains a separate editorial note. This check does not claim to re-review every assertion of the Examples chapter.

\subsubsection{Integration scope}\label{integration-scope}

The preparation records must retain each original statement, the four exact source operations, both original chapter files, the prior complete cumulative chapter files, source/proof locators, and the receiving Algebra group IDs. Removing these four edits must recover the complete prior chapters, including their existing additions. The changes must also produce standalone fixes-only patches against the original authority, with no unrelated theorem additions.

This source-ordered propagation does not make the Algebra R58 candidate or either cross-chapter correction public or admitted. Those transitions require their own source replay, build evidence, rendered checks, and publication receipts.
